\section{问题描述}

给定足够光滑且有界的关节期望轨迹 \(q_d(t),\dot q_d(t),\ddot q_d(t)\)。定义位置误差、速度误差和误差状态为
\[
e=q_m-q_d,\qquad \dot e=\dot q_m-\dot q_d,
\qquad z=\begin{bmatrix}e^\top&\dot e^\top\end{bmatrix}^\top\in\mathbb R^{2n}.
\]
本文首先研究关节空间轨迹跟踪问题，控制目标是设计关节力矩 \(\tau\)，使误差在给定时间上界后满足
\begin{equation}
e(t)=0,\qquad \dot e(t)=0,\qquad t\ge T.
\label{eq:tracking_goal}
\end{equation}
其中，\(T\) 由观测阶段和 PDF 阶段的设计参数共同确定。需要区分的是，在理想连续时间、精确模型和精确扰动补偿条件下，式 \eqref{eq:tracking_goal} 表示预设时间精确跟踪；当存在模型误差、持续未补偿扰动、执行器饱和或数值离散时，本文只主张预设时间附近的实际有界跟踪或小邻域收敛。

本文采用如下假设。

\textbf{假设1：} 在所考虑工作区域内，\(M_e(q_m)\) 一致正定且可逆，\(M_e(q_m)\) 与 \(C_e(q_m,\dot q_m)\dot q_m\) 可由完整模型获得，或由名义模型近似获得。

\textbf{假设2：} 参考轨迹 \(q_d,\dot q_d,\ddot q_d\) 有界且连续，关节位置和速度可测，且参考信号在控制器启动时具有一致的历史定义。

\textbf{假设3：} 人工延迟项 \(z(t-h)\) 可由历史缓存获得，其中 \(h>0\) 为设计延迟。延迟初值函数在 \([-h,0]\) 上有界且分段连续。

\textbf{假设4：} 对于预设时间扰动观测器，等效加速度扰动
\[
\Delta_a(t)=M_e^{-1}(q_m)d(t)
\]
及其一阶导数有界；其上界可以未知。

\textbf{假设5：} 严格的预设时间精确跟踪结论只在连续时间实现、无执行器饱和、模型补偿精确以及观测器按理论条件实现时成立。存在观测误差、模型不确定性、采样和量化误差时，相关项统一视为残余扰动进行实际有界性分析。

上述假设明确了本文理论结果的适用范围。后文先在理想条件下建立 PDF 闭环的预设时间精确跟踪结论，再讨论残余扰动对实际性能的影响。末端任务空间轨迹若由广义雅可比映射生成关节参考，则属于关节空间控制器的上层参考生成问题，本文的主要理论对象仍为式 \eqref{eq:tracking_goal} 所示的关节空间误差系统。
